\chapter{Group-Theoretic Foundations of Pitch, Tuning, and Musical Harmony}
\label{ch:group-theory}

Every procedural synthesizer must map discrete musical notation onto continuous acoustic frequencies, amplitudes, and spatial positions. In B-Sound, this mapping is not an ad-hoc collection of constants; it is grounded in the algebraic structure of pitch space, cyclic and dihedral symmetry groups, and constant-power Riemannian geometry on the stereo panning arc.

\section{Logarithmic Pitch Space as a One-Dimensional Lie Group}

Human auditory perception of pitch is fundamentally logarithmic (the Weber--Fechner law for frequency). Two tones with frequencies $f_1, f_2 \in \mathbb{R}_{>0}$ are perceived as separated by the same musical interval whenever their ratio $f_2 / f_1$ is constant.

\begin{definition}[Pitch-Frequency Isomorphism]
The positive real numbers under multiplication $(\mathbb{R}_{>0}, \times)$ form a one-dimensional connected abelian Lie group. Choosing a reference tonic $F_0 > 0$ (in B-Sound's analytical synthesizer, middle $\mathrm{C}_4 = 261.6256\,\text{Hz}$; in ISO~16 concert pitch, $\mathrm{A}_4 = 440.00\,\text{Hz}$), the base-2 logarithm defines a Lie group isomorphism:
\begin{equation}
\label{eq:log2-iso}
\Phi : (\mathbb{R}_{>0}, \times) \xrightarrow{\;\sim\;} (\mathbb{R}, +), \qquad y = \Phi(f) = \log_2\!\left(\frac{f}{F_0}\right), \qquad f(y) = F_0 \cdot 2^y.
\end{equation}
\end{definition}

Under $\Phi$, multiplicative frequency ratios become additive pitch intervals:
\[
\Phi(f_1 \cdot r) = \Phi(f_1) + \log_2(r).
\]
In particular, doubling a frequency (an \emph{octave}, $r = 2$) corresponds to integer translation $y \mapsto y + 1$.

\section{Octave Equivalence and the Cyclic Group $\mathbb{Z}_{12}$}

Because frequencies separated by powers of $2$ share harmonic overtone alignment, Western and B-Sound music theory identify pitches modulo the discrete subgroup $\mathbb{Z} \subset (\mathbb{R}, +)$:
\[
\mathbb{T}^1 = \mathbb{R} / \mathbb{Z} \cong S^1.
\]
Discretizing the pitch-class circle $\mathbb{T}^1$ into $N$ equal subdivisions yields the $N$-tone equal temperament ($N$-TET) cyclic group $\mathbb{Z}_N = \mathbb{Z} / N\mathbb{Z}$. Setting $N = 12$ yields standard 12-tone equal temperament (12-TET), whose generator is the semitone ratio $2^{1/12} \approx 1.0594631$.

\begin{theorem}[Optimality of $N = 12$ for Fifth and Third Approximation]
In the continued-fraction expansion of $\log_2(3/2) \approx 0.5849625$, the convergents are $\frac{1}{2}, \frac{3}{5}, \frac{7}{12}, \frac{24}{41}, \frac{31}{53}, \dots$. The convergent $7/12$ is the smallest denominator exceeding $5$ that simultaneously approximates both the pure fifth ($3/2$, error $-1.955$ cents) and the pure fourth ($4/3$, error $+1.955$ cents) within $2$ cents while keeping a compact keyboard alphabet.
\end{theorem}

\begin{proof}
Compute the cent deviation $\Delta(k, N) = 1200 \left(\frac{k}{N} - \log_2(3/2)\right)$. For $N = 12, k = 7$, $\Delta(7, 12) = 700 - 701.955 = -1.955\,\text{cents}$. Since $\log_2(4/3) = 1 - \log_2(3/2)$, the complementary step $5/12$ approximates $4/3$ with error $+1.955\,\text{cents}$. The next sharper convergent requires $N = 41$ or $N = 53$, multiplying instrument polyphony tables by fourfold.
\end{proof}

In B-Sound and \emph{Orchestra}, every MIDI note $n \in \{0, 1, \dots, 127\}$ maps to Hz via Equation~(5) of the specification:
\begin{equation}
\label{eq:midi-freq}
f(n) = 440.0 \times 2^{(n - 69)/12}.
\end{equation}
To keep the real-time audio path free of transcendental \texttt{powf()} calls, B-Sound precomputes the 128-entry table \texttt{s\_pitch[128]} once at boot. Every incoming note traverses a strict four-stage algebraic pipeline before table lookup:
\begin{enumerate}[leftmargin=1.5em]
  \item \textbf{Key-split filter:} discard $n \notin [\texttt{key\_lo}, \texttt{key\_hi}]$ (where $(0,0)$ denotes $[0,127]$; \texttt{orch\_cmd\_split} cycles $[0,127]$, $[0,59]$, $[60,127]$ modulo $3$).
  \item \textbf{Track transposition:} $n \leftarrow n + \Delta_{\mathrm{trk}}$.
  \item \textbf{Program transposition:} $n \leftarrow n + \Delta_{\mathrm{prog}}$ (e.g., Viola $\Delta = -7$, Contrabass $\Delta = -12$).
  \item \textbf{Domain clamp and lookup:} $n' = \mathrm{clamp}(n, 0, 127)$, $f = \texttt{s\_pitch}[n']$.
\end{enumerate}

\section{Transposition, Inversion, and the Dihedral Group $D_{12}$}

Musical transformations on pitch classes $x \in \mathbb{Z}_{12}$ are generated by two fundamental symmetries:
\begin{definition}[Transposition and Inversion Operators]
For any $k \in \mathbb{Z}_{12}$, define the \emph{transposition} $T_k$ and \emph{inversion} $I_k$ on $\mathbb{Z}_{12}$ by
\begin{equation}
T_k(x) \equiv x + k \pmod{12}, \qquad I_k(x) \equiv -x + k \pmod{12}.
\end{equation}
\end{definition}

\begin{proposition}[Dihedral Structure of Pitch-Class Symmetries]
The set $\mathcal{T}/\mathcal{I} = \{T_k, I_k : k \in \mathbb{Z}_{12}\}$ under function composition forms a group isomorphic to the dihedral group $D_{12}$ of order $24$, with semidirect product decomposition $D_{12} \cong \mathbb{Z}_{12} \rtimes \mathbb{Z}_2$ and multiplication laws:
\[
T_a \circ T_b = T_{a+b}, \qquad T_a \circ I_b = I_{a+b}, \qquad I_a \circ T_b = I_{a-b}, \qquad I_a \circ I_b = T_{a-b}.
\]
\end{proposition}

\begin{proof}
Direct evaluation on $x \in \mathbb{Z}_{12}$: $(I_a \circ T_b)(x) = -(x + b) + a = -x + (a - b) = I_{a-b}(x)$, and $(I_a \circ I_b)(x) = -(-x + b) + a = x + (a - b) = T_{a-b}(x)$. Every $I_k$ is an involution ($I_k^2 = T_0 = \mathrm{id}$), and $\{T_k\} \trianglelefteq D_{12}$ is a normal subgroup of index $2$.
\end{proof}

Major and minor triads are single orbits under $D_{12}$: the C-major triad $\{0, 4, 7\}$ maps under $I_7$ to $\{7, 3, 0\} = \{0, 3, 7\}$, the C-minor triad. Thus inversion exchanges major and minor mode while preserving interval-class content $\langle 0, 0, 1, 1, 1, 0 \rangle$.

\section{Neo-Riemannian $PLR$ Group and Smooth Voice-Leading}

In late-Romantic and cinematic harmony (including the \emph{Analytical Etudes} and \emph{House of the Rising Sun} in B-Sound's repertoire), progressions frequently move between triads by minimal semitone displacement rather than strict functional roots. Let $\mathcal{S}$ denote the set of $24$ major and minor triads in $\mathbb{Z}_{12}$.

\begin{definition}[Neo-Riemannian Involutions $\mathcal{P}, \mathcal{L}, \mathcal{R}$]
On any major triad $\langle r, r+4, r+7 \rangle$ and minor triad $\langle r, r+3, r+7 \rangle$ (arithmetic in $\mathbb{Z}_{12}$), define three parsimonious voice-leading involutions that hold two common tones fixed and move the third by one or two semitones:
\begin{align}
\mathcal{P}(\text{Major } r) &= \text{Minor } r & (\text{moves } r+4 \leftrightarrow r+3,\; 1\text{ semitone}), \\
\mathcal{L}(\text{Major } r) &= \text{Minor } (r+4) & (\text{moves } r \leftrightarrow r-1,\; 1\text{ semitone}), \\
\mathcal{R}(\text{Major } r) &= \text{Minor } (r+9) & (\text{moves } r+7 \leftrightarrow r+9,\; 2\text{ semitones}).
\end{align}
\end{definition}

\begin{theorem}[Simply Transitive Dual Action of $PLR$]
The group $\mathcal{G}_{PLR} = \langle \mathcal{P}, \mathcal{L}, \mathcal{R} \rangle$ is isomorphic to $D_{12}$ of order $24$, commutes with the $\mathcal{T}/\mathcal{I}$ group ($\mathcal{G}_{PLR}$ is the centralizer of $D_{12}$ in $\mathrm{Sym}(\mathcal{S})$), and acts simply transitively on the $24$ consonant triads.
\end{theorem}

\begin{example}[Analytical Application to B-Sound Repertoire]
Consider the iconic A-minor progression of \emph{The House of the Rising Sun} (Track~3) and the \emph{Analytical Etudes} (`Am -> C -> D -> F -> Am -> E`):
\begin{itemize}[leftmargin=1.5em]
  \item $\mathrm{Am} \xrightarrow{\;\mathcal{R}\;} \mathrm{C}$ holds $\{C, E\}$ fixed and moves $A \to G$.
  \item $\mathrm{C} \xrightarrow{\;\mathcal{R}\mathcal{L}\;} \mathrm{Am}$ and $\mathrm{Am} \xrightarrow{\;\mathcal{L}\;} \mathrm{F}$ holds $\{A, C\}$ fixed and moves $E \to F$.
  \item $\mathrm{Am} \xrightarrow{\;\mathcal{P}\;} \mathrm{A}$ and $\mathrm{C} \xrightarrow{\;\mathcal{L}\mathcal{P}\;} \mathrm{E}$ illustrate why small $PLR$ word lengths produce smooth, low-energy trajectories on the analytical pitch canvas $y(x)$.
\end{itemize}
\end{example}

\section{Velocity Curves, Channel Gain, and Constant-Power Panning}

In B-Sound, dynamics and spatialization are governed by three closed-form laws (Equations~6--8 of \emph{Orchestra}). Let $u = v_m / 127 \in [0, 1]$ for MIDI velocity $v_m \in \{0, \dots, 127\}$. The linear family of three velocity curves maps $u$ into $[0.08, 0.45]$:
\begin{equation}
\label{eq:vel-curves}
v_0(u) = 0.08 + 0.37\,u, \qquad v_1(u) = 0.08 + 0.37\,\sqrt{u}, \qquad v_2(u) = 0.08 + 0.37\,u^2.
\end{equation}
Classical dynamic markings ($pp$: $0.08$--$0.12$, $mp$: $0.14$--$0.18$, $mf$: $0.20$--$0.26$, $f$: $0.28$--$0.36$, $ff$: $0.38$--$0.45$) are named bands of this single scale rather than a second hidden curve. Track fader gain (Equation~7) obeys:
\begin{equation}
\label{eq:track-gain}
G_{\mathrm{trk}} =
\begin{cases}
0, & \text{if } \texttt{volume} = 0 \text{ or } \texttt{muted} = 1, \\
0.08 + 0.37 \cdot (\texttt{volume} / 127), & \text{otherwise.}
\end{cases}
\end{equation}

Stereo panning maps MIDI CC10 onto the unit circle $S^1$ via the constant-power ($-3\,\text{dB}$ center) law (Equation~8):
\begin{equation}
\label{eq:pan-law}
p = \frac{\mathrm{CC10} - 64}{64} \in \left[-1, \frac{63}{64}\right], \qquad
G_L(p) = \cos\!\left(\frac{\pi}{4}(p + 1)\right), \qquad
G_R(p) = \sin\!\left(\frac{\pi}{4}(p + 1)\right).
\end{equation}
Because $G_L(p)^2 + G_R(p)^2 = 1$ identically for all $p$, acoustic power remains invariant as a sound source sweeps across the stereo stage, with $G_L(0) = G_R(0) = 1/\sqrt{2} \approx 0.7071$ at center.
